\begin{lemma}{6.3.1}\label{I.6.3.1}
Let \(f\colon\mathbf{Y}\to\mathbf{X}\) be a \(G\)-equivariant morphism and \(\mathbf{M}\) a
\(G\)-equivariant \(\mathbf{O}_{\mathbf{X}}\)-module. Then the inverse image
\(f^{*}(\mathbf{M})=\mathbf{Y}\times_{f}\mathbf{M}\) is a \(G\)-equivariant
\(\mathbf{O}_{\mathbf{Y}}\)-module.
\end{lemma}

On the other hand, \(\mathbf{G}\) acts on
\(\Hom_{\widehat{\mathcal{C}}}(\mathbf{Y},\mathbf{X})\). Indeed, let
\(T\in\Ob(\mathcal{C})\), \(g\in\mathbf{G}(T)\) and \(\varphi\) a \(T\)-morphism
\(\mathbf{Y}_{T}\to\mathbf{X}_{T}\); then \(g\) defines automorphisms
\(\lambda_{\mathbf{Y}}(g)\) and \(\lambda_{\mathbf{X}}(g)\) of \(\mathbf{Y}_{T}\) and
\(\mathbf{X}_{T}\), and one will denote by \(g\cdot\varphi\) (or also
\(g\varphi g^{-1}\)) the morphism
\(\lambda_{\mathbf{X}}(g)\circ\varphi\circ\lambda_{\mathbf{Y}}(g^{-1})\).
This defines an action of \(\mathbf{G}(T)\) on
\(\Hom_{T}(\mathbf{Y}_{T},\mathbf{X}_{T})\), functorial in \(T\).

\begin{definition}{6.3.2}\label{I.6.3.2}
If \(\varphi\colon\mathbf{Y}\to\mathbf{X}\) is an arbitrary
\(\widehat{\mathcal{C}}\)-morphism, one may therefore consider its stabilizer
\(\operatorname{Stab}_{\mathbf{G}}(\varphi)\) (\cf 2.3.3.1): for every
\(T\in\Ob(\mathcal{C})\), \(\operatorname{Stab}_{\mathbf{G}}(\varphi)(T)\) is the
subgroup of \(\mathbf{G}(T)\) formed of the \(g\) such that
\(g\circ\varphi_{T}=\varphi_{T}\circ g\), \ie such that the diagram
\[
\xymatrix{
\mathbf{Y}_{T} \ar[r]^{\varphi_{T}} \ar[d]_{g}
  & \mathbf{X}_{T} \ar[d]^{g} \\
\mathbf{Y}_{T} \ar[r]_{\varphi_{T}}
  & \mathbf{X}_{T}
}
\]
commutes. Then the morphism \(\varphi\colon\mathbf{Y}\to\mathbf{X}\) is
equivariant under \(\operatorname{Stab}_{\mathbf{G}}(\varphi)\).
\end{definition}

\par\addvspace{0.55\baselineskip}
\noindent\textbf{6.4. Global sections.}\ ---\ \label{I.sec.6.4}
Let \(\mathbf{M}\) be a \(G\)-equivariant \(\mathbf{O}_{\mathbf{X}}\)-module.
Denote by \(S_{0}\) the final object of \(\mathcal{C}\) and (\cf Exp.~II, 1.1)
by \(\prod_{\mathbf{X}/S_{0}}\mathbf{M}\) the ``functor of sections of
\(\mathbf{M}\) over \(\mathbf{X}\)": it is the functor which to every
\(T\in\Ob(\mathcal{C})\) associates
\[
\Hom_{\mathbf{X}}(\mathbf{X}_{T},\mathbf{M})
  =\Hom_{\mathbf{X}_{T}}(\mathbf{X}_{T},\mathbf{M}_{T})
  =\Gamma(\mathbf{M}_{T}/\mathbf{X}_{T}).
\]
Let us recall on the other hand that every morphism
\(g\colon\mathbf{Z}\to\mathbf{Y}\) of \(\widehat{\mathcal{C}}\)-objects over
\(\mathbf{X}\) induces a morphism of abelian groups
\(\mathbf{M}(g)\colon\mathbf{M}(\mathbf{Y})\to\mathbf{M}(\mathbf{Z})\), which is
compatible with the ring morphism
\(g^{*}\colon\mathbf{O}(\mathbf{Y})\to\mathbf{O}(\mathbf{Z})\). In particular,
when \(\mathbf{Z}=\mathbf{Y}\) (\(g\) then being an
\(\mathbf{X}\)-endomorphism of \(\mathbf{Y}\)), one obtains a morphism of
abelian groups
\[
\mathbf{M}(g)\colon\mathbf{M}(\mathbf{Y})\longrightarrow\mathbf{M}(\mathbf{Y})
\]
which is not in general a morphism of \(\mathbf{O}(\mathbf{Y})\)-modules, but
which satisfies, for every \(m\in\mathbf{M}(\mathbf{Y})\) and
\(\alpha\in\mathbf{O}(\mathbf{Y})\):
\[
\mathbf{M}(g)(\alpha\cdot m)=g^{*}(\alpha)\cdot\mathbf{M}(g)(m).
\]
This being said, one will simply write, in what follows, \(g^{*}\) instead of
\(\mathbf{M}(g)\).

Let \(T\in\Ob(\mathcal{C})\) and let \(\mathbf{X}_{T}=\mathbf{X}\times T\) and
\(\mathrm{pr}_{\mathbf{X}}\) be the projection
\(\mathbf{X}_{T}\to\mathbf{X}\). For every
\(\alpha\in\mathbf{O}(\mathbf{X}_{T})\) and \(g\in\mathbf{G}(T)\), set
\(g(\alpha)=(g^{-1})^{*}(\alpha)\): this is the element of
\(\mathbf{O}(\mathbf{X}_{T})\) defined set-theoretically by: for every
\(S\in\Ob(\mathcal{C})\) and \(x\in\mathbf{X}(S)\), \(t\in T(S)\),
\[
g(\alpha)(x,t)=\alpha(g^{-1}x,t).
\]
One thus obtains an action (on the left) of \(\mathbf{G}(T)\) by ring
automorphisms on \(\mathbf{O}(\mathbf{X}_{T})\), functorial in \(T\), and such
that \(g(\alpha)=\alpha\) if \(\alpha\) is the image in
\(\mathbf{O}(\mathbf{X}_{T})\) of an element of \(\mathbf{O}(T)\).

Let us now denote by \(\phi\) the identity morphism of \(\mathbf{X}_{T}\)
(\cf 6.3 and the generalization below to a morphism
\(\phi\colon\mathbf{Y}\to\mathbf{X}\)) and denote
\(\Hom_{\mathbf{X}}(\mathbf{X}_{T},\mathbf{M})\) by
\(\mathbf{M}(\phi g^{-1})\) \resp \(\mathbf{M}(\phi)\), according as
\(\mathbf{X}_{T}\) is regarded as an \(\mathbf{X}\)-object via
\(\mathrm{pr}_{\mathbf{X}}\circ\lambda(g^{-1})_{T}\), \resp
\(\mathrm{pr}_{\mathbf{X}}\).

Then \(\lambda(g^{-1})_{T}\) is an \(\mathbf{X}\)-morphism between these two
\(\mathbf{X}\)-objects, hence, by what precedes, one obtains a morphism of
abelian groups
\[
(g^{-1})^{*}\colon\mathbf{M}(\phi)\longrightarrow\mathbf{M}(\phi g^{-1}),
  \qquad m\mapsto m\circ\lambda(g^{-1})_{T}
\]
which satisfies \((g^{-1})^{*}(\alpha\cdot m)=(g\alpha)\cdot(g^{-1})^{*}(m)\)
for every \(m\in\mathbf{M}(\phi)\) and \(\alpha\in\mathbf{O}(\mathbf{X}_{T})\).
(If \(m\) is a section of \(\mathbf{M}_{T}\) over \(\mathbf{X}_{T}\), then
\((g^{-1})^{*}(m)\) is the section defined set-theoretically by
\((x,t)\mapsto m(g^{-1}x,t)\).) In particular, \((g^{-1})^{*}\) is a morphism
of \(\mathbf{O}(T)\)-modules.

By the functoriality of the morphisms of \(\mathbf{O}(\mathbf{X}_{T})\)-modules
\(\Lambda_{x}(g)\), one obtains a commutative diagram:
\[
\xymatrix{
\mathbf{M}(\phi) \ar[r]^{(g^{-1})^{*}} \ar[d]_{\Lambda_{\phi}(g)}
  & \mathbf{M}(\phi g^{-1}) \ar[d]^{\Lambda_{\phi g^{-1}}(g)} \\
\mathbf{M}(g\phi) \ar[r]_{(g^{-1})^{*}}
  & \mathbf{M}(g\phi g^{-1})
}
\]
and \(g\phi g^{-1}=\phi\), since \(\phi\) is the identity map. Writing
\(A(g)=(g^{-1})^{*}\circ\Lambda_{\phi}(g)\), one therefore obtains a
morphism of abelian groups
\[
A(g)\colon\mathbf{M}(\phi)\longrightarrow\mathbf{M}(\phi)
\]
which is ``compatible with the action of \(\mathbf{G}(T)\) on
\(\mathbf{O}(\mathbf{X}_{T})\)'', \ie which satisfies
\[
A(g)(\alpha\cdot m)=(g\alpha)\cdot A(g)(m).
\]
Finally, if \(h\) is a second element of \(\mathbf{G}(T)\), it follows from
the functoriality of \(\mathbf{M}\) and of the morphisms \(\Lambda_{x}(g)\)
that one has a commutative diagram
\[
\xymatrix@C=1.6em@R=1.8em{
& \mathbf{M}(\phi) \ar[dl]_{\Lambda_{\phi}(g)} \ar[dr]^{\Lambda_{\phi}(hg)} & & \\
\mathbf{M}(g\phi) \ar[rr]^{\Lambda_{\phi}(h)} \ar[d]_{(g^{-1})^{*}}
  && \mathbf{M}(hg\phi) \ar[d]^{(g^{-1})^{*}} \ar[dr]^{((hg)^{-1})^{*}} & \\
\mathbf{M}(\phi) \ar[rr]_{\Lambda_{\phi}(h)}
  && \mathbf{M}(h\phi) \ar[r]_{(h^{-1})^{*}}
  & \mathbf{M}(\phi)
}
\]
whence \(A(h)\circ A(g)=A(hg)\). Consequently, one
has obtained the following proposition.

\begin{proposition}{6.4.1}\label{I.6.4.1}
For every \(T\), the \(\mathbf{O}(\mathbf{X}_{T})\)-module
\(\bigl(\prod_{\mathbf{X}/S_{0}}\mathbf{M}\bigr)(T)=\Gamma(\mathbf{M}_{T}/\mathbf{X}_{T})\)
is endowed, in a manner functorial in \(T\), with an action of
\(\mathbf{G}(T)\) ``compatible with the action of \(\mathbf{G}(T)\) on
\(\mathbf{O}(\mathbf{X}_{T})\)'', \ie which satisfies
\[
A(g)(\alpha\cdot m)=g(\alpha)\cdot A(g)(m).
\]
As \(g(\alpha)=\alpha\) for \(\alpha\in\mathbf{O}(T)\), this shows, in
particular, that \(\prod_{\mathbf{X}/S_{0}}\mathbf{M}\) is a
\(\mathbf{G}\text{-}\mathbf{O}_{S_{0}}\)-module.
\end{proposition}

More generally let, as in 6.3, \(\mathbf{Y}\) be a second \(G\)-equivariant
\(\widehat{\mathcal{C}}\)-object, \(\phi\colon\mathbf{Y}\to\mathbf{X}\) a
\(\widehat{\mathcal{C}}\)-morphism and
\(\mathbf{H}=\operatorname{Stab}_{\mathbf{G}}(\phi)\). Then the fiber product
\(\mathbf{M}_{\phi}=\mathbf{Y}\times_{\phi}\mathbf{M}\) is an
\(\mathbf{H}\)-equivariant \(\mathbf{O}_{\mathbf{Y}}\)-module. One therefore
obtains the:

\begin{corollary}{6.4.2}\label{I.6.4.2}
The functor \(\prod_{\mathbf{Y}/S_{0}}\mathbf{M}_{\phi}\) is a
\(\operatorname{Stab}_{\mathbf{G}}(\phi)\text{-}\mathbf{O}_{S_{0}}\)-module.
\end{corollary}

\par\addvspace{0.55\baselineskip}
\noindent\textbf{6.5. \(G\)-equivariant \(\OX\)-modules.}\ ---\ \label{I.sec.6.5}
Let us apply what precedes in the following situation. Let \(S\) be a scheme,
\(G\) an \(S\)-group scheme acting on an \(S\)-scheme \(X\), and \(\mathcal{F}\)
an \(\OX\)-module (not necessarily quasi-coherent).

\begin{definition}{6.5.1}\label{I.6.5.1}
One says that \(\mathcal{F}\) is a \(G\)-equivariant \(\OX\)-module if the
\(\mathbf{O}_{X}\)-module \(\mathbf{M}=\mathbf{W}(\mathcal{F})\) is
\(G\)-equivariant, \ie if one has given oneself, for every morphism
\((g,x)\colon U\to G\times_{S}X\), isomorphisms of \(\mathbf{O}_{U}\)-modules
\[
\Lambda_{x}(g)\colon\mathbf{W}\bigl(x^{*}(\mathcal{F})\bigr)
  \isomto\mathbf{W}\bigl((gx)^{*}(\mathcal{F})\bigr),
\]
functorial in \(U\) and satisfying
\(\Lambda_{hx}(g)\circ\Lambda_{x}(h)=\Lambda_{x}(gh)\). As the functor
\(\mathbf{W}\) is fully faithful (\cf 4.6.1.1 and 4.6.2), one therefore
obtains isomorphisms of \(\mathcal{O}_{U}\)-modules
\[
\Lambda_{x}(g)\colon x^{*}(\mathcal{F})\isomto(gx)^{*}(\mathcal{F}),
\]
where one recalls that \(gx\) denotes the morphism
\(\lambda\circ(g,x)\colon U\to X\). In particular, applying this to the
identity morphism of \(G\times_{S}X\), one obtains an isomorphism of
\(\mathcal{O}_{G\times_{S}X}\)-modules
\begin{equation*}
(\star)\qquad
\theta\colon\mathrm{pr}_{X}^{*}(\mathcal{F})\isomto\lambda^{*}(\mathcal{F})
\end{equation*}
such that the diagram \((\star\star)\) of morphisms of sheaves on
\(G\times_{S}G\times_{S}X\) below is commutative:
\[
\xymatrix@C=0.7em@R=2.6em{
*+{\scriptstyle\mathrm{pr}_{2,3}^{*}\circ\mathrm{pr}_{X}^{*}(\mathcal{F})}
  \ar[r]^{\scriptstyle\mathrm{pr}_{2,3}^{*}(\theta)}_{\sim}
&
*+{\scriptstyle\mathrm{pr}_{2,3}^{*}\circ\lambda^{*}(\mathcal{F})}
&
*+{\scriptstyle(\mathrm{id}_{G}\times\lambda)^{*}\circ\mathrm{pr}_{X}^{*}(\mathcal{F})}
  \ar[d]^{\scriptstyle(\mathrm{id}_{G}\times\lambda)^{*}(\theta)}
\\
*+{\scriptstyle(\mu\times\mathrm{id}_{X})^{*}\circ\mathrm{pr}_{X}^{*}(\mathcal{F})}
  \ar[r]^{\scriptstyle(\mu\times\mathrm{id}_{X})^{*}(\theta)}_{\sim}
&
*+{\scriptstyle(\mu\times\mathrm{id}_{X})^{*}\circ\lambda^{*}(\mathcal{F})}
  \ar@{}[u]|{(\star\star)}
&
*+{\scriptstyle(\mathrm{id}_{G}\times\lambda)^{*}\circ\lambda^{*}(\mathcal{F})}\,.
}
\]
Moreover, if \(\mathcal{E}\) is a second \(G\)-equivariant \(\OX\)-module, one
says that a morphism of \(\OX\)-modules
\(\phi\colon\mathcal{E}\to\mathcal{F}\) is \(G\)-equivariant if the morphism
\(\mathbf{W}(\phi)\colon\mathbf{W}(\mathcal{E})\to\mathbf{W}(\mathcal{F})\) is.
With the notations above, this amounts to saying that
\(\theta_{\mathcal{F}}\circ\mathrm{pr}_{X}^{*}(\phi)
  =\lambda^{*}(\phi)\circ\theta_{\mathcal{E}}\).
\end{definition}

\begin{remark}{6.5.2}\label{I.6.5.2}
If \(f\colon Y\to X\) is a \(G\)-equivariant morphism, it follows from 6.3.1
that \(f^{*}(\mathcal{F})\) is a \(G\)-equivariant \(\mathcal{O}_{Y}\)-module.
\end{remark}

\begin{remark}{6.5.3}\label{I.6.5.3}
Suppose moreover that \(\mathcal{F}\) is \emph{quasi-coherent}. Then it
follows from what precedes that the datum of a structure of \(G\)-equivariant
\(\mathbf{O}_{X}\)-module on \(\mathbf{M}=\mathbf{W}(\mathcal{F})\)
\emph{amounts} to the datum of such a structure on \(\mathbf{V}(\mathcal{F})\),
which still \emph{amounts} to giving oneself an action of \(G\) on the vector
fibration \(\mathbf{V}(\mathcal{F})\), compatible with the action on \(X\) and
``linear'' on the fibers.

Indeed, denote by \(\phi\) the isomorphism
\(\lambda^{*}(\mathcal{F})\isomto\mathrm{pr}_{X}^{*}(\mathcal{F})\) inverse to
\(\theta\). For every morphism \((g,x)\colon U\to G\times_{S}X\), one has an
isomorphism of \(\mathcal{O}_{U}\)-modules
\(\phi_{x}(g)=\Lambda_{x}(g)^{-1}=\Lambda_{gx}(g^{-1})\) from
\((gx)^{*}(\mathcal{F})\) to \(x^{*}(\mathcal{F})\); it induces an isomorphism
of \(\mathbf{O}_{U}\)-modules
\[
\mathbf{V}\bigl((gx)^{*}(\mathcal{F})\bigr)
  \isomto\mathbf{V}\bigl(x^{*}(\mathcal{F})\bigr)
\]
which one will denote by \({}^{t}\phi_{x}(g)\). One then has
\[
{}^{t}\phi_{gx}(h)\circ{}^{t}\phi_{x}(g)
  ={}^{t}\bigl(\Lambda_{gx}(h)\circ\phi_{x}(g)\bigr)^{-1}
  ={}^{t}\Lambda_{x}(hg)^{-1}
  ={}^{t}\phi_{x}(hg).
\]
As, for every \(X\)-scheme \(f\colon Y\to X\), one has
\(\mathbf{V}(\mathcal{F})\times_{f}Y\simeq\mathbf{V}(f^{*}(\mathcal{F}))\)
(\cf 4.6.2), one therefore obtains that the isomorphism
\begin{align*}
{}^{t}\phi\colon
  G\times_{S}\mathbf{V}(\mathcal{F})
  &=(G\times_{S}X)\times_{\mathrm{pr}_{X}}\mathbf{V}(\mathcal{F})
   =\mathbf{V}\bigl(\mathrm{pr}_{X}^{*}(\mathcal{F})\bigr) \\
  &\isomto\mathbf{V}\bigl(\lambda^{*}(\mathcal{F})\bigr)
   =(G\times_{S}X)\times_{\lambda}\mathbf{V}(\mathcal{F})
\end{align*}
endows \(\mathbf{V}(\mathcal{F})\) with a structure of \(G\)-equivariant
\(\mathbf{O}_{X}\)-module. Finally, identifying each functor
\(\mathbf{V}(-)\) with the vector fibration which represents it, and composing
\({}^{t}\phi\) with the projection onto \(\mathbf{V}(\mathcal{F})\), one
obtains an action of \(G\) on the vector fibration \(\mathbf{V}(\mathcal{F})\),
compatible with the action on \(X\) and ``linear'' on the fibers.

One thus recovers the definition given, for example, in [GIT, Chap.~1,
\S~3], except that in \loccit Mumford considers an \(\OX\)-module
\(\mathcal{E}\) locally free of finite rank, and defines an action of \(G\)
on \(\mathbf{V}(\mathcal{E})=\mathbf{W}(\mathcal{E}^{\vee})\). Indeed, the
preceding diagram \((\star\star)\), with \(\theta\) replaced by \(\phi\) and
the direction of the arrows reversed, is exactly the one found in \loccit,
Def.~1.6, and the isomorphism \({}^{t}\phi\) above coincides with the
isomorphism \(\Phi\) of \loccit, p.~31.
\end{remark}

\begin{remark}{6.5.4}\label{I.6.5.4}
Let us consider in particular the case where \(X=S\), endowed with the
trivial action of \(G\). In this case, a \(G\)-equivariant \(\OS\)-module
\(\mathcal{F}\) is the same thing as a \(G\text{-}\mathbf{O}_{S}\)-module
(\cf 4.7.1). Moreover, if one denotes by \(f\) the morphism \(G\to S\) (here
equal both to \(\mathrm{pr}_{X}\) and to \(\lambda\)), then the isomorphism
\begin{equation*}
(\star)\qquad
\theta\colon f^{*}(\mathcal{F})\isomto f^{*}(\mathcal{F})
\end{equation*}
is an element of
\begin{align*}
\Hom_{\mathcal{O}_{G}}\bigl(f^{*}(\mathcal{F}),f^{*}(\mathcal{F})\bigr)
  &=\Hom_{\mathbf{O}_{G}}\bigl(\mathbf{W}(f^{*}(\mathcal{F})),
      \mathbf{W}(f^{*}(\mathcal{F}))\bigr) \\
  &=\End_{\mathbf{O}_{S}}\bigl(\mathbf{W}(\mathcal{F})\bigr)(G)
\end{align*}
which is none other than the morphism
\(\rho\colon G\to\End_{\mathbf{O}_{S}}(\mathbf{W}(\mathcal{F}))\) defining the
operation of \(G\) on \(\mathbf{W}(\mathcal{F})\). Moreover, \(\theta\)
corresponds by adjunction to the morphism of \(\OS\)-modules
\(f_{*}(\theta)\circ\tau\colon\mathcal{F}\to f_{*}f^{*}(\mathcal{F})\), where
\(\tau\colon\mathcal{F}\to f_{*}f^{*}(\mathcal{F})\) is the ``unit'' morphism
of the adjunction. (This will be used in
\(\mathrm{VI}_{\mathrm{B}}\), 11.10.bis.)
\end{remark}

\par\addvspace{0.55\baselineskip}
\noindent\textbf{6.6. The functors \(\prod_{X/S}\mathbf{W}(\mathcal{F})\) and
\(\mathbf{W}(p_{*}(\mathcal{F}))\).}\ ---\ \label{I.sec.6.6}
Let \(S\) be a scheme, \(G\) an \(S\)-group scheme acting on \(S\)-schemes \(X\)
and \(Y\) via the morphisms \(\lambda\colon G\times_{S}X\to X\) and
\(\mu\colon G\times_{S}Y\to Y\), and let \(p\colon X\to Y\) be a
\(G\)-equivariant morphism. Suppose that the morphisms \(p\colon X\to Y\) and
\(\nu_{Y}\colon Y\to S\) are \emph{quasi-compact} and \emph{quasi-separated},
and that \(\pi\colon G\to S\) is \emph{flat}.

Then the projection \(\mathrm{pr}_{Y}\colon G\times_{S}Y\to Y\) is flat, as is
\(\mu\) (since \(\mu\) is the composite of \(\mathrm{pr}_{Y}\) and of the
automorphism \((g,y)\mapsto(g,gy)\)). Finally, for a variable \(S\)-scheme
\(f\colon T\to S\), one will denote by \(p^{T}\colon X_{T}\to T\) and
\(f_{X}\colon X_{T}\to X\) the morphisms deduced from \(p\) and \(f\) by base
change.
% typo?: printed ``f : T → variable''; read as f : T → S variable

Let \(\mathcal{F}\) be a \emph{quasi-coherent} and \emph{\(G\)-equivariant}
\(\OX\)-module,
and let \(\theta\) be the isomorphism
\(\mathrm{pr}_{X}^{*}(\mathcal{F})\isomto\lambda^{*}(\mathcal{F})\) of
6.5.1\((\star)\). As \(p\) is quasi-compact and quasi-separated,
\(p_{*}(\mathcal{F})\) is quasi-coherent (\cf EGA~I, 9.2.1).

\begin{lemma}{6.6.1}\label{I.6.6.1}
The quasi-coherent \(\mathcal{O}_{Y}\)-module \(p_{*}(\mathcal{F})\) is
\(G\)-equivariant.
\end{lemma}

Indeed, one has the two cartesian squares below:
\[
\xymatrix{
G\times_{S}X \ar[r]^{\mathrm{pr}_{X}} \ar[d]_{q}
  & X \ar[d]^{p}
  & G\times_{S}X \ar[l]_{\lambda} \ar[d]^{q} \\
G\times_{S}Y \ar[r]_{\mathrm{pr}_{Y}}
  & Y
  & G\times_{S}Y \ar[l]^{\mu}\,.
}
\]
As \(p\) is quasi-compact and quasi-separated and \(\mathrm{pr}_{Y}\) and
\(\mu\) are flat, it follows from EGA~III, 1.4.15 (supplemented by
EGA~IV\(_{1}\), 1.7.21) that one has
\(q^{*}\mathrm{pr}_{X}^{*}(\mathcal{F})=\mathrm{pr}_{Y}^{*}p_{*}(\mathcal{F})\)
and \(q^{*}\lambda^{*}(\mathcal{F})=\mu^{*}p_{*}(\mathcal{F})\). Consequently,
\(\theta\) induces an isomorphism
\[
\theta_{Y}\colon\mathrm{pr}_{Y}^{*}p_{*}(\mathcal{F})
  \isomto\mu^{*}p_{*}(\mathcal{F}),
\]
and one likewise obtains that \(\theta_{Y}\) satisfies the ``cocycle
condition'' 6.5.1\((\star\star)\) (since the morphisms
\(G\times_{S}G\times_{S}Y\to G\times_{S}Y\) which occur are flat). This proves
the lemma.

In particular, take \(Y=S\) endowed with the trivial action of \(G\). Taking
remark 6.5.4 into account, one then obtains that \(p_{*}(\mathcal{F})\) is a
\(G\text{-}\mathbf{O}_{S}\)-module. If moreover \(G\) is affine over \(S\) and
if one writes \(\mathcal{A}(G)=\pi_{*}(\mathcal{O}_{G})\), then
\(p_{*}(\mathcal{F})\) is therefore an \(\mathcal{A}(G)\)-comodule, by 4.7.2.

On the other hand, by 6.4.1, the functor
\(\prod_{X/S}\mathbf{W}(\mathcal{F})\), which to every \(f\colon T\to S\)
associates
\[
\mathbf{W}\bigl(f_{X}^{*}(\mathcal{F})\bigr)(X_{T})
  =\Gamma\bigl(X_{T},f_{X}^{*}(\mathcal{F})\bigr)
  =\Gamma\bigl(T,p^{T}_{*}f_{X}^{*}(\mathcal{F})\bigr)
\]
is a \(G\text{-}\mathbf{O}_{S}\)-module. Moreover, one has a canonical
morphism
\(\tau\colon\mathbf{W}(p_{*}(\mathcal{F}))\to\prod_{X/S}\mathbf{W}(\mathcal{F})\),
which is given for every \(f\colon T\to S\) by the canonical morphism
\[
\Gamma\bigl(T,f^{*}p_{*}(\mathcal{F})\bigr)
  \longrightarrow\Gamma\bigl(T,p^{T}_{*}f_{X}^{*}(\mathcal{F})\bigr)
\]
and which is an isomorphism when one restricts it to the full subcategory of
schemes \(T\) flat over \(S\), and one checks without difficulty that \(\tau\)
is a morphism of \(G\text{-}\mathbf{O}_{S}\)-modules. Consequently, one has
obtained the following proposition (for point~(ii), compare with [GIT],
p.~32).

\begin{proposition}{6.6.2}\label{I.6.6.2}
Let \(S\) be a scheme, \(G\) an \(S\)-group scheme acting on an \(S\)-scheme
\(X\), and \(\mathcal{F}\) a quasi-coherent and \(G\)-equivariant
\(\OX\)-module. Suppose that \(\pi\colon G\to S\) is flat and that
\(p\colon X\to S\) is quasi-compact and quasi-separated.
\par
\begin{enumeratei}
\item Then \(p_{*}(\mathcal{F})\) is a quasi-coherent
\(G\text{-}\mathbf{O}_{S}\)-module. On the other hand, the canonical morphism
\(\mathbf{W}(p_{*}(\mathcal{F}))\to\prod_{X/S}\mathbf{W}(\mathcal{F})\) is a
morphism of \(G\text{-}\mathbf{O}_{S}\)-modules, and these two functors
coincide on the category of flat \(S\)-schemes.
\item If moreover \(G\) is affine over \(S\) and if one writes
\(\mathcal{A}(G)=\pi_{*}(\mathcal{O}_{G})\), then \(p_{*}(\mathcal{F})\) is
endowed with a structure of \(\mathcal{A}(G)\)-comodule.%
{\renewcommand{\thefootnote}{(55)}\nde{The same holds if \(G\) is flat,
quasi-compact and quasi-separated over \(S\) and if \(p_{*}(\mathcal{F})\) is
a flat \(\OS\)-module, \cf Exp.~\(\mathrm{VI}_{\mathrm{B}}\),
11.6.1\,(ii).}}
\end{enumeratei}
\end{proposition}

\par\addvspace{0.55\baselineskip}
\noindent\textbf{6.7. Stabilizers.}\ ---\ \label{I.sec.6.7}
Let \(S\) be a scheme, \(G\) an \(S\)-group scheme acting on an \(S\)-scheme
\(X\), and let \(\mathcal{F}\) be a quasi-coherent \(G\)-equivariant
\(\OX\)-module. Let \(Y\) be a second \(S\)-scheme endowed with an action of
\(G\) (possibly trivial), let \(\phi\colon Y\to X\) be an \(S\)-morphism, not
necessarily \(G\)-equivariant, and let \(\operatorname{Stab}_{G}(\phi)\) be
the stabilizer in \(G\) of \(\phi\) (\cf 6.3.2).

Let us note at once (see Exp.~\(\mathrm{VI}_{\mathrm{B}}\), \S~6) that
\(\operatorname{Stab}_{G}(\phi)\) is representable by a closed group
subscheme \(H\) of \(G\) if \(X\) is \emph{separated} over \(S\) and if \(Y\) is
\emph{essentially free} over \(S\) (\cf \loccit, Def.~6.2.1). Indeed, consider
the morphism
\(r\colon G\times_{S}Y\to X\times_{S}X\) given set-theoretically by
\(r(g,y)=\bigl(\phi(y),\,g\phi(g^{-1}y)\bigr)\), and let \(P=G\times_{S}Y\) and
\(P'\) be the inverse image by \(r\) of the diagonal \(\Delta_{X/S}\). Then one
has (\cf \loccit, 6.2.4\,(a))
\[
\operatorname{Stab}_{G}(\phi)=\prod_{P/G}P'
\]
and hence, by \loccit, \(\operatorname{Stab}_{G}(\phi)\) is representable by a
closed group subscheme \(H\) of \(G\) if \(X\) is separated over \(S\) and if
\(Y\) is essentially free over \(S\); this second condition being
automatically satisfied if \(S\) is the spectrum of a field, or else if
\(Y=S\). Under these hypotheses, \(\phi^{*}(\mathcal{F})\) is then a
quasi-coherent \(H\)-equivariant \(\mathcal{O}_{Y}\)-module (\cf 6.5.2).

Therefore, if moreover \(\pi\colon G\to S\) and \(p\colon Y\to S\) are
quasi-compact and quasi-separated over \(S\), and \(\pi\) is flat, then
\(p_{*}\phi^{*}(\mathcal{F})\) is an \(H\text{-}\mathbf{O}_{S}\)-module, by
6.6.2. In particular, one obtains the:

\begin{corollary}{6.7.1}\label{I.6.7.1}
Let \(S\) be a scheme, \(G\) an \(S\)-group scheme acting on an \(S\)-scheme
\(X\), and \(\mathcal{F}\) a quasi-coherent \(G\)-equivariant \(\OX\)-module.
Suppose that \(\pi\colon G\to S\) is flat and that \(X\to S\) is quasi-compact
and separated. Let \(\tau\colon S\to X\) be a section of \(X\) over \(S\).
\par
\begin{enumeratei}
\item The stabilizer \(H=\operatorname{Stab}_{G}(\tau)\) is a closed group
subscheme of \(G\), and \(\tau^{*}(\mathcal{F})\) is a quasi-coherent
\(H\text{-}\mathbf{O}_{S}\)-module.
\item If moreover \(H\) is affine over \(S\) (for example, if \(G\) is) and if
one writes \(\mathcal{A}(H)=\pi_{*}(\mathcal{O}_{H})\), then
\(\tau^{*}(\mathcal{F})\) is an \(\mathcal{A}(H)\)-comodule.
\end{enumeratei}
\end{corollary}

\par\addvspace{0.55\baselineskip}
\noindent\textbf{6.8. \(G\)-equivariant sheaves on \(G\).}\ ---\ \label{I.sec.6.8}
To conclude, let us indicate two results (6.8.1 and 6.8.6 below) which will
be used in exposés~II and~III (\cf in particular III, 4.25).

\begin{proposition}{6.8.1}\label{I.6.8.1}
Let \(S\) be a scheme, \(\pi\colon G\to S\) an \(S\)-group scheme, and
\(\varepsilon\colon S\to G\) the unit section. One considers the action by
left translations of \(G\) on itself. Then the functors
\(\mathcal{E}\mapsto\pi^{*}(\mathcal{E})\) and
\(\mathcal{F}\mapsto\varepsilon^{*}(\mathcal{F})\) induce equivalences,
quasi-inverse to one another, between the category of quasi-coherent
\(\OS\)-modules and that of quasi-coherent \(G\)-equivariant
\(\mathcal{O}_{G}\)-modules.
\end{proposition}

\begin{proof}
Denote by \(\mu\) the multiplication of \(G\) and by \(\mathrm{pr}_{2}\) the
second projection \(G\times_{S}G\to G\). As \(\pi\circ\mu=\pi\circ\mathrm{pr}_{2}\),
for every quasi-coherent \(\OS\)-module \(\mathcal{E}\) one has a canonical
isomorphism \(\mu^{*}\pi^{*}(\mathcal{E})=\mathrm{pr}_{2}^{*}\pi^{*}(\mathcal{E})\),
and one easily checks that this isomorphism satisfies the ``cocycle
condition'' 6.5.1\((\star\star)\), \ie \(\pi^{*}(\mathcal{E})\) is a
\(G\)-equivariant \(\mathcal{O}_{G}\)-module. As
\(\varepsilon^{*}\pi^{*}(\mathcal{E})=\mathcal{E}\), the functor
\(\mathcal{E}\mapsto\pi^{*}(\mathcal{E})\) is fully faithful; it therefore
remains to see that for every \(G\)-equivariant quasi-coherent
\(\mathcal{O}_{G}\)-module \(\mathcal{F}\), one has
\(\mathcal{F}\simeq\pi^{*}\varepsilon^{*}(\mathcal{F})\).
% proof continues on the next page; tombstone suppressed at the chunk boundary
\renewcommand{\qedsymbol}{}
\end{proof}
